\section{A rank certificate for the Hankel seed}
\label{app:certificate}

We prove Lemma~\ref{lem:seed}.  Recall that
\[
 H=\R[x_0,\ldots,x_7]_2=L\oplus P,
 \qquad \dim H=36,\quad \dim L=8,\quad \dim P=28,
\]
where $L$ consists of the quadratic monomials divisible by $x_0$.
We seek a positive semidefinite Hankel matrix $T$ of rank $20$ whose
restriction to $L$ is positive definite.  Its projected kernel
$K=\operatorname{proj}_P(\ker T)$ must have independent pairwise
products and span all pure cubics after multiplication by pure linear
forms.  The remaining requirement is
\[
 \dim\bigl((\ker T)\cdot H_0\bigr)=321,
 \qquad H_0=\{p\in H:T(p,L)=0\},
\]
inside the $330$-dimensional quartic space.

Our starting matrix is the moment matrix of twenty points arranged in
two cyclic orbits.  At that matrix the last product space has dimension
$302$.  We exhibit a tangent direction for which nineteen independent
product directions appear to first order.  All nonvanishing assertions
are certified below by explicit minors over $\mathbb F_7$; we then
explain why these certificates produce a real positive semidefinite
matrix with the required properties.

\subsection{Cyclic orbit matrices and the required rank increase}
\label{subsec:orbit-seed}

For the algebraic construction work first over $\mathbb C$, keeping
$x_0$ as the homogenizing coordinate.  Give the eight variables the
respective weights
\begin{equation}
 \omega=(0,5,1,9,2,8,4,6)\pmod {10}.
 \label{eq:certificate-weights}
\end{equation}
Let $\zeta$ be a primitive tenth root of unity.  For
$a=(1,a_1,\ldots,a_7)$, take the two orbits of $(1,\ldots,1)$ and $a$
under the diagonal action
$x_j\mapsto\zeta^{\omega_j}x_j$.
Give every orbit point weight $1/10$, and let $T(a)$ be the resulting
quadratic moment matrix.  If $m$ is a quartic monomial, its moment is
\begin{equation}
 t_a(m)=
 \begin{cases}
  1+m(a),&\operatorname{wt}(m)=0,\\
  0,&\operatorname{wt}(m)\ne0.
 \end{cases}
 \label{eq:orbit-moments}
\end{equation}
In particular $T(a)_{u,v}=t_a(uv)$ for quadratic monomials $u,v$.
Formula~\eqref{eq:orbit-moments} has integer coefficients in the
parameters and also defines the matrices used in the finite-field
calculation; no tenth roots of unity in $\mathbb F_7$ are needed.

Write $I_i$ for the ordered quadratic monomials of weight $i$ and set
\[
 C_i(a)=
 \begin{pmatrix}
  1&\cdots&1\\
  u_1(a)&\cdots&u_{|I_i|}(a)
 \end{pmatrix},
 \qquad (u_1,\ldots,u_{|I_i|})=I_i.
\]
The only nonzero block in column weight $i$ of $T(a)$ is
$C_{-i}(a)^TC_i(a)$.  If every $C_i(a)$ has rank two, then
$\rank T(a)=20$ and its kernel is the direct sum of the nullspaces
$J_i=\ker C_i(a)$.  Thus $J=\ker T(a)$ has dimension $16$.

Let $D$ denote a matrix whose columns are the quartic products of a
basis of $J$ with a basis of $H_0$.  For a Hankel variation $T_1$
satisfying
\begin{equation}
 T_1|_{J\times J}=0,
 \label{eq:certificate-tangent}
\end{equation}
write $D_1$ for the first-order product variation obtained by
differentiating the equations for bases of $J$ and $H_0$.  We give
explicit formulas below and later realize $T_1$ as the velocity of
a smooth arc.  We shall establish, simultaneously with the two
product conditions on $K$,
\begin{equation}
 \rank D=302,
 \qquad
 \rank\begin{pmatrix}D&0\\D_1&D\end{pmatrix}\ge623
 =2\cdot302+19.
 \label{eq:certificate-target-ranks}
\end{equation}
These ranks do not depend on the choices of differentiable bases.
Indeed, a change of the derivatives of those bases changes $D_1$ by
a matrix of the form $AD+DB$, which is removed by block row and column
operations in the second matrix of
\eqref{eq:certificate-target-ranks}.

We first record the upper bounds that give the target ranks their
meaning.  They apply in characteristic zero wherever $T$ has rank
$20$, $J\cap L=0$, and $K$ spans all pure cubics after multiplication
by pure linear forms.  The seven translations and the dilation in
the chart $x_0=1$ have independent derivatives on the kernel graph,
by the proof of Lemma~\ref{lem:kernel-affine-action}, which uses
only these hypotheses and works over $\mathbb C$ as well as over
$\R$.  Recall the obstruction to a dependence involving dilation:
it would allow a translation making the kernel equal to the homogeneous space $K$.
Cubic spanning would then annihilate all chart moments of degrees
three and four, forcing rank at most $8+1=9$ instead of $20$.

A translation or dilation preserves the projected kernel $K$.
Continue each kernel vector with its highest homogeneous part
fixed, so that its first derivative lies in $L$.  Differentiating
the kernel equation therefore gives
$T'(0)J\subset T(L)$, so its quartic functional annihilates $JH_0$.
Scaling $T$ gives one further independent annihilator, since scaling
has zero kernel derivative whereas the eight affine motions have
independent kernel derivatives.  Hence, throughout the rank-$20$
locus under consideration,
\begin{equation}
 \rank D\le330-9=321.
 \label{eq:certificate-upper-321}
\end{equation}

At an orbit matrix, scaling belongs to the twenty-dimensional family
obtained by varying the individual weights of its twenty points.
Rank $20$ means
that their evaluation functionals are independent already on
quadratics, and hence on quartics by multiplication by $x_0^2$.
Each such quartic evaluation annihilates $JH_0$, and all weight
variations have zero kernel derivative.  The eight affine motions
remain independent modulo these twenty directions.  Thus at the
orbit matrix
\begin{equation}
 \rank D\le330-20-8=302.
 \label{eq:certificate-upper-302}
\end{equation}
In particular, attaining rank at least $302$ there is an open condition
equivalent to equality.

\subsection{Bases and first derivatives modulo seven}
\label{subsec:certificate-bases}

All entries in the remainder of this subsection and the next are in
$\mathbb F_7$.  Use the specialization
\begin{equation}
 a=(1,6,6,0,3,0,2,5).
 \label{eq:certificate-point}
\end{equation}
Index monomials by nondecreasing strings of variable indices, ordered
lexicographically, both globally and within each weight.  For example,
$23$ means $x_2x_3$ and $0011$ means $x_0^2x_1^2$.
Every row and column number below starts at one within its indicated
block.

For $C_i$, take pivot columns
\begin{equation}
 E_i=(1,2)\quad\text{except that}\quad E_0=E_8=(1,3).
 \label{eq:certificate-pivots}
\end{equation}
Use the nullspace basis having an identity matrix on the remaining
coordinates, in their order.  The quadratic blocks and resulting columns
of $J_i$ are as follows.  A digit string in the third column records
one entire column vector in the order of $I_i$.
\begin{center}
\begin{tabular}{c l l}
\toprule
$i$&$I_i$&columns of $J_i$\\
\midrule
0&$00,11,23,45,67$&$61000,\ 00610,\ 40201$\\
1&$02,17,34$&$421$\\
2&$04,22,56,77$&$4210,\ 2401$\\
3&$15,24,36$&$601$\\
4&$06,13,44,57$&$6010,\ 0601$\\
5&$01,26,37$&$511$\\
6&$07,12,46,55$&$4210,\ 2401$\\
7&$14,27,35$&$151$\\
8&$05,33,47,66$&$6100,\ 3031$\\
9&$03,16,25$&$601$\\
\bottomrule
\end{tabular}
\end{center}
Each specified two-column pivot is invertible.  Thus every $C_i$ has
rank two and $T$ has rank $20$.  The block on $L$ is also invertible:
in the order $00,01,\ldots,07$ its determinant is $6$.

For later calculations, the ordered list of weight-zero quartics is
\begin{equation}
\begin{gathered}
0000,0011,0023,0045,0067,0126,0137,0225,0334,0447,\\
0466,0556,0577,1111,1123,1145,1167,1244,1257,1346,\\
1355,2233,2247,2266,2345,2367,3356,3377,4446,4455,\\
4567,4777,5557,5666,6677.
\end{gathered}
\label{eq:certificate-zero-quartics}
\end{equation}
Define $T_1$ by setting its quartic moments outside weight zero equal
to zero and its moments in the order
\eqref{eq:certificate-zero-quartics} equal to
\begin{equation}
 (2,0,0,0,3,6,1,4,1,6,0,3,0,5,2,2,0,0,4,0,6,
 \underbrace{0,\ldots,0}_{14}).
 \label{eq:certificate-tangent-moments}
\end{equation}
The weight-zero block of products in $\Sym^2J$ has seventeen columns.
Its transpose times the vector
\eqref{eq:certificate-tangent-moments} is the zero vector.  Products
of every other weight pair trivially with $T_1$, so
\eqref{eq:certificate-tangent} holds.

We specify a basis $N_i$ for the weight-$i$ part of $H_0$.  There is
one coordinate of $L$ in each weight except $3$ and $7$.  Denote its
within-block index by $L_i$; thus $L_0=1$ corresponds to $00$.
If weight $-i$
has no $L$ coordinate, let $N_i$ be the identity matrix.  Otherwise
let $N_i$ be the nullspace basis of the single row $T_{L_{-i},I_i}$,
using pivot $L_i$ and identity on the free indices.  Its pivot is
nonzero because $T|_L$ is invertible.

The following formulas specify derivatives of both bases.  Put
$J_i'$ equal to zero off the pivot rows $E_i$, and set
\begin{equation}
 (J_i')_{E_i}=
 -\bigl(T_{E_{-i},E_i}\bigr)^{-1}
       (T_1)_{E_{-i},I_i}J_i.
 \label{eq:certificate-J-derivative}
\end{equation}
If $N_i$ is an identity matrix, put $N_i'=0$.  In the other cases,
put $N_i'$ equal to zero off row $L_i$, and set
\begin{equation}
 (N_i')_{L_i}=
 -\bigl(T_{L_{-i},L_i}\bigr)^{-1}
       (T_1)_{L_{-i},I_i}N_i.
 \label{eq:certificate-N-derivative}
\end{equation}
These are the differentiated kernel equations with the free
coordinates held fixed.  For $J_i'$, tangency ensures that solving
the two selected equations also solves all the equations
$TJ_i'+T_1J_i=0$: the image of $T_1J_i$ lies in the image of $T$,
and the two selected rows give coordinates on that image in the
relevant weight block.  The formulas for $N_i'$ follow by
differentiating $T(L,N_i)=0$.

\subsection{The explicit minors}
\label{subsec:certificate-minors}

We now give the minors that establish the two product conditions on
$K$ and the rank increase in
\eqref{eq:certificate-target-ranks}.  All product matrices use the
following convention.  In weight $w$, multiply a column $u$ from
block $i$ by a column $v$ from block $w-i$, adding
\[
 u_{jj'}v_{kk'}
 \quad\text{to row}\quad \operatorname{sort}(jj'kk').
\]
Order columns lexicographically by the weight $i$, the column number
of $u$, and the column number of $v$.  For products from
$\Sym^2J$, retain only pairs for which the first index
$(i,\text{column number})$ is at most the second.  Thus each
unordered product occurs once, with its ordinary polynomial
coefficients.

Let $G_1$ be this matrix for $J\cdot J$, restricted to quartic rows
with no zero index.  It is the product matrix for $\Sym^2K$.
Let $G_2$ be the product matrix of $J$ with the seven coordinate
vectors $x_0x_j$, $1\le j\le7$, restricted to rows with exactly one
zero index.  Removing that zero identifies it with the product map
from $K$ times pure linear forms to pure cubics.  Rows and columns
are numbered anew after the restriction to a weight and to these
row sets.  In all the tables, a colon denotes the inclusive range.

The following square minors use all columns of $G_1$ and all rows of
$G_2$, respectively.  The specified omissions determine the other
index set completely.
\begin{center}
\small
\begin{tabular}{c l r r l r r}
\toprule
&\multicolumn{3}{c}{$G_1$}&\multicolumn{3}{c}{$G_2$}\\
$w$&omit rows&size&det.&omit columns&size&det.\\
\midrule
0&$15,19{:}22$&17&3&$8,10,11$&8&6\\
1&$9,11,12,15{:}20$&11&4&$9{:}11$&8&6\\
2&$13,15,19{:}22$&16&4&$8,10,11$&9&2\\
3&$3,10,14{:}20$&11&5&$8,10$&8&6\\
4&$13,15,18,19,21,22$&16&6&$10{:}12$&9&6\\
5&$9,11,13,15{:}20$&11&5&$8,10,11$&8&4\\
\bottomrule
\end{tabular}
\end{center}
The number of columns of $G_1$ in weights $0,\ldots,5$ is
$17,11,16,11,16,11$; the number of rows of $G_2$ is
$8,8,9,8,9,8$.  Thus these determinants give precisely the full
column and full row ranks needed in these weights.

For $D$, use first factors from $J$ and second factors from $H_0$,
with the preceding product order.  Form $D_1$ by replacing the first
factor by its derivative in
\eqref{eq:certificate-J-derivative}, replacing the second by its
derivative in \eqref{eq:certificate-N-derivative}, and adding the
two resulting product columns.  In weights $0,\ldots,5$, the
dimensions of $D$ are respectively
\[
 35\times48,\quad31\times42,\quad35\times47,\quad
 31\times43,\quad35\times47,\quad31\times42.
\]
The next table specifies a square submatrix $E=D[R,S]$ in each
weight by listing the omitted rows and columns; $R,S$ always
denote the retained indices.
\begin{center}
\small
\setlength{\tabcolsep}{4pt}
\begin{tabular}{c l l r r}
\toprule
$w$&omit rows&omit columns&size&$\det E$\\
\midrule
0&$13,34,35$&$5,9,11,33,34,36{:}38,40{:}43,45{:}48$&32&6\\
1&$24,26,31$&$7,9,10,28,30,32,33,36{:}42$&28&1\\
2&$32,34,35$&$12,14{:}16,18,19,35,38,40,41,43{:}47$&32&6\\
3&$26,31$&$14,16,17,19,20,35{:}43$&29&3\\
4&$33,34,35$&$17,20,21,23{:}25,27,28,41{:}47$&32&5\\
5&$19,26,31$&$17,18,20,22,23,25,26,36{:}42$&28&3\\
\bottomrule
\end{tabular}
\end{center}

To obtain the bordered minors, use the indices $R,S$ in both
diagonal positions of $\left(\begin{smallmatrix}D&0\\D_1&D
\end{smallmatrix}\right)$.  Add rows $R'$ in its second block row
and columns $S'$ in its first block column.  Set
$A=D[R,S']$ and $B=D[R',S]$.  Eliminating the two copies of $E$
leaves the corner
\begin{equation}
\begin{split}
 Z={}&D_1[R',S']-D_1[R',S]E^{-1}A
       -BE^{-1}D_1[R,S']\\
    &\hspace{27mm}+BE^{-1}D_1[R,S]E^{-1}A.
\end{split}
\label{eq:certificate-corner}
\end{equation}
The corner matrices are the following, with rows separated by
semicolons when their entries are listed.
\begin{center}
\begin{tabular}{c l l l r}
\toprule
$w$&$R'$&$S'$&$Z$&$\det Z$\\
\midrule
0&$34$&$38$&$(5)$&5\\
1&$24,31$&$38,40$&$(4,2;\ 2,4)$&5\\
2&$34,35$&$45,46$&$(6,4;\ 2,0)$&6\\
3&$26,31$&$37,39$&$(3,6;\ 0,5)$&1\\
4&$33,35$&$43,46$&$(6,4;\ 1,2)$&1\\
5&$26,31$&$37,38$&$(2,6;\ 2,0)$&2\\
\bottomrule
\end{tabular}
\end{center}
In particular the resulting bordered minor has size $2|R|+1$ in
weight zero and $2|R|+2$ in each of weights one through five.
The bordered determinant is, up to its row and column ordering
sign, $(\det E)^2\det Z$, as follows directly by eliminating the
two diagonal copies of $E$.  Thus the table proves the asserted
bordered lower bounds without requiring an upper bound on the
rank of $D$ in the finite field.

For completeness, every entry in these tables can be checked with
the following fixed elimination rule.  At column $n$, interchange
row $n$ with the first row at or below it whose entry $t$ in that
column is nonzero.  Divide the new row by $t$ and subtract the
appropriate multiples from every other row, always reducing modulo
seven.  The determinant is the product of the pivots, with one
minus sign for each interchange.  Appended columns give the
inverse actions in \eqref{eq:certificate-corner} by the same rule.
For the $G_2$ minor in weight zero, the signed pivots are
\[
 (-2,-2,-5,-1,4,-1,1,5).
\]
For $E$ in weight zero they are
\begin{equation}
\begin{gathered}
 1,4,1,4,-4,-1,-4,-5,-5,-4,-4,6,-4,4,6,-3,\\
 -3,-1,6,6,4,5,-5,2,1,5,2,5,6,4,4,4.
\end{gathered}
\label{eq:certificate-pivot-example}
\end{equation}
Together with the monomial order, the basis rules, and the moment
vectors above, this specifies the finite calculation using matrices
with at most thirty-five quartic rows in each weight.

\subsection{Characteristic zero and simultaneous real data}
\label{subsec:certificate-lifting}

We now pass from the explicit modular minors to actual orbit
matrices and tangent directions in characteristic zero.  First lift
\eqref{eq:certificate-point} to the same integer vector.  The chosen
pivots of the $C_i$, the entries used as pivots for $N_i$, and the
$2\times2$ blocks in
\eqref{eq:certificate-J-derivative} are all invertible modulo seven.
Hence the corresponding bases and inverse formulas are defined over
$\mathbb Z_{(7)}$, the rational numbers whose denominators are not
divisible by seven.

There is one point to address before lifting the bordered minors:
the first variation must remain tangent exactly, not just modulo
seven.  The weight-zero part of $\Sym^2J$ has dimension seventeen.
Its seventeen product columns are independent, since their
restriction to the pure rows contains the nonzero $17\times17$
minor of $G_1$ above.  Consequently tangency of a weight-zero
functional gives seventeen independent equations on its thirty-five
moment coordinates.  Choose as pivot coordinates the seventeen
quartic rows of that minor.  Lift the remaining eighteen coordinates
of \eqref{eq:certificate-tangent-moments} to integers and solve
these seventeen equations over $\mathbb Z_{(7)}$.  The solution
reduces to precisely the displayed moment vector.  It therefore
defines an exact characteristic-zero tangent whose basis
derivatives reduce to
\eqref{eq:certificate-J-derivative} and
\eqref{eq:certificate-N-derivative}.  Every nonzero determinant in
the preceding tables thus proves that its defining rational
function is nonzero in characteristic zero.

It remains to obtain all weights and a real positive semidefinite
matrix simultaneously.  We spell out the genericity argument to
make clear that the tangent need not be chosen separately for
different weights.  The permutation
\[
 (x_2,x_3),\quad(x_4,x_5),\quad(x_6,x_7)
 \quad\text{interchanged, with }x_0,x_1\text{ fixed},
\]
negates every weight and preserves the family of orbit matrices.
The nonzero minors for weights $1,2,3,4$ therefore also give nonempty
rank conditions for weights $9,8,7,6$, respectively.

More formally, take the Zariski open subset of the parameter space
$\mathbb C^7$ on which all $C_i$ have rank two, $T|_L$ is
invertible, and the seventeen weight-zero products are independent.
This is a nonempty irreducible parameter space.  The solutions of
the seventeen tangency equations form a rank-eighteen algebraic
vector bundle over it: locally an invertible seventeen-column
minor solves for seventeen moment coordinates in terms of the
other eighteen.  Its total space is irreducible.  Each condition
certified above is a nonempty Zariski open subset of this space,
as is its image under the weight-negating permutation, after
restricting to the common nonempty parameter domain.  Their finite
intersection is nonempty.  The same reasoning for the minors of
$G_1$ and $G_2$, which depend only on $a$, gives both product
conditions on $K$.

The resulting lower bounds in weights $0,\ldots,9$ are
\begin{equation}
\begin{array}{c|rrrrrrrrrr|r}
w&0&1&2&3&4&5&6&7&8&9&\text{sum}\\\hline
\rank G_1&17&11&16&11&16&11&16&11&16&11&136\\
\rank G_2&8&8&9&8&9&8&9&8&9&8&84\\
\rank D&32&28&32&29&32&28&32&29&32&28&302\\
\rank\left(\begin{smallmatrix}D&0\\D_1&D\end{smallmatrix}\right)
 &65&58&66&60&66&58&66&60&66&58&623
\end{array}
\label{eq:certificate-rank-sums}
\end{equation}
Here $136=\dim\Sym^2K$ and $84$ is the dimension of the pure
cubic space.  The first two rows prove the required injectivity
and surjectivity.  Equation~\eqref{eq:certificate-upper-302}
then turns the third row into equality, and the fourth row gives
\eqref{eq:certificate-target-ranks}.

Impose now the real structure in which $x_0,x_1$ are real and the
pairs $(x_2,x_3)$, $(x_4,x_5)$, $(x_6,x_7)$ are conjugate.  Its
fixed parameter set
\[
 a_1\in\R,\qquad a_3=\overline{a_2},\quad
 a_5=\overline{a_4},\quad a_7=\overline{a_6}
\]
is a real form of $\mathbb C^7$ and is Zariski dense.  One way to
see this is to write each conjugate pair as $(u+iv,u-iv)$: this is
an invertible complex linear change of coordinates from a real
coordinate space, and a polynomial vanishing on that real space
is zero.  The weight-zero tangent equations at a fixed real
parameter are preserved by the real structure.  Their solution
space is therefore the complexification of its real fixed
subspace, which is likewise Zariski dense.  A nonempty open
condition in the tangent bundle is attained over a nonempty open
set of parameters, and within each such fiber it is an open
condition on a linear space.  We can thus choose simultaneously
a parameter in the real form and a real tangent satisfying every
rank condition in \eqref{eq:certificate-rank-sums}.

For these parameters all twenty orbit points are real in the
corresponding real coordinates.  Indeed multiplication by
$\zeta^{\omega_j}$ respects the conjugate pairs and multiplies
the weight-five coordinate by $\pm1$.  Since every point has
positive weight $1/10$, the moment matrix is real positive
semidefinite.  Its rank is $20$, and its nonsingular restriction
to $L$ is positive definite.
Choose ordinary real coordinates on this fixed real form, keeping
$x_0$ and the pure-variable span, and relabel them $x_0,\ldots,x_7$.
The induced polynomial change of coordinates preserves $L\oplus P$
and multiplication, and hence the total product and first-variation
ranks above. Thus these are real Hankel data in the conventions of
Section~\ref{sec:hankel}.

\subsection{Deformation to the seed matrix}

We have obtained a real positive semidefinite orbit matrix $T$ and
a real Hankel tangent $T_1$ satisfying
\eqref{eq:certificate-target-ranks}.  Independence of the products
in $\Sym^2K$ implies independence of those in $\Sym^2J$ by taking
highest chart degree.  Lemma~\ref{lem:hankel-stratum}, or its
Schur-complement proof, therefore shows that the rank-$20$ Hankel
locus is smooth near $T$, with tangent condition
\eqref{eq:certificate-tangent}.  There is a real smooth arc
$T(t)$ on this locus with $T(0)=T$ and $T'(0)=T_1$.
For sufficiently small $t$, its twenty nonzero eigenvalues remain
positive, its restriction to $L$ remains positive definite, and
the two product conditions on $K(t)$ remain valid.

Choose smooth local kernel and $H_0$ bases along the arc and let
$D(t)$ be their product matrix.  Apply fixed invertible row and
column operations that put $D(0)$ in rank normal form
\[
 D(0)=\begin{pmatrix}I_{302}&0\\0&0\end{pmatrix}.
\]
Writing the transformed derivative in the same blocks as
$\left(\begin{smallmatrix}A&B\\C&F\end{smallmatrix}\right)$,
elementary block elimination gives
\[
 \rank\begin{pmatrix}D(0)&0\\D'(0)&D(0)\end{pmatrix}
 =604+\rank F.
\]
Thus $\rank F\ge19$.  Eliminate the invertible leading
$302\times302$ block of $D(t)$.  Its remaining Schur complement
is $tF+o(t)$, since both off-diagonal blocks vanish at $t=0$.
A nonzero $19\times19$ minor of $F$ remains nonzero after division
by $t^{19}$ for every sufficiently small nonzero $t$.  Hence
$\rank D(t)\ge321$ there.  The upper bound
\eqref{eq:certificate-upper-321} applies along the arc and gives
equality.  Taking one such $T(t)$ proves Lemma~\ref{lem:seed}.

\paragraph{Reproducible calculation.}
The accompanying source file \texttt{verify\_certificate.py}
reconstructs the monomial blocks, bases, differentiated products,
and every displayed minor using exact integer arithmetic modulo
seven.  It also verifies the ranks in weights six through nine
at the very same specialization
\eqref{eq:certificate-point} and tangent
\eqref{eq:certificate-tangent-moments}.  This supplies an
independent reproduction of the finite calculation; the explicit
minors and the characteristic-zero argument above establish the
lemma without relying on the program.
